\origpage{87}{N. Robertson, P. Seymour / Journal of Combinatorial Theory, Series B 89 (2003) 77--108}
\noindent $(A, B) \in \mathscr{L}_h$ satisfies $A \subseteq C_h$ and $D_h \subseteq B$. If in addition $w \ge 0$ and $\mathscr{L}_h \cup \{(D_h, C_h)\}$ has tree-width $\le w$ for $1 \le h \le k$, we say that $\mathscr{L}^*$ is an \textit{enlargement of $\mathscr{L}$ by tree-width $\le w$.}

4.3. \textit{Let $\lambda$ be a tie-breaker in a hypergraph $G$, let $\theta \ge 2$, and let $\mathscr{T}_1, \dots, \mathscr{T}_n$ be distinct tangles in $G$, each of order $\theta$, where $n \ge 1$. For $1 \le i \le n$ let $\mathscr{L}_i \subseteq \mathscr{T}_i$ be a location of order $\le \frac{3}{4}\theta$ which $\theta$-isolates $\mathscr{T}_i$; and suppose that for every tangle $\mathscr{T}$ in $G$ of order $\theta$, there is a unique $i$ with $1 \le i \le n$ such that $\mathscr{L}_i \subseteq \mathscr{T}$. Let $(T, \tau)$ be a standard tree-decomposition of $G$ relative to $\mathscr{T}_1, \dots, \mathscr{T}_n$, where $V(T) = \{t_1, \dots, t_n\}$ and $t_i$ represents $\mathscr{T}_i$ for $1 \le i \le n$. Then for $1 \le i \le n$, the location of $t_i$ in $(T, \tau)$ is an enlargement of $\mathscr{L}_i$ by tree-width $\le \frac{9}{4}\theta$.}

\textbf{Proof.} Let $1 \le i \le n$, and let $\mathscr{L}_i = \{(C_1, D_1), \dots, (C_k, D_k)\}$. Let $\mathscr{L}^*$ be the location of $t_i$ in $(T, \tau)$.

\hspace*{12pt}(1) \textit{If $(A, B) \in \mathscr{L}^*$ then $(A, B)$ has order $\le \frac{3}{4}\theta$ and there exists $h$ with $1 \le h \le k$ such that $A \subseteq C_h$ and $D_h \subseteq B$.}

\textbf{Subproof.} Since $(A, B) \in \mathscr{L}^*$, there exists $j \ne i$ with $1 \le j \le n$ such that $t_i, t_j$ are adjacent in $T$ and $(A, B)$ is the $(\mathscr{T}_i, \mathscr{T}_j)$-distinction. Since $\mathscr{L}_j \subseteq \mathscr{T}_j$ and $j \ne i$ it follows that $\mathscr{L}_{\delm{\backslash}\fix{i}} \delm{\subseteq}\fix{\not\subseteq} \mathscr{T}_j$,\fixnote{fix:4.3-claim1}{The slash of the ``not'' sign has slipped into the subscript, and $\mathscr{L}_i \not\subseteq \mathscr{T}_j$ is meant.} and so there exists $h$ with $1 \le h \le k$ such that $(C_h, D_h) \notin \mathscr{T}_j$. Since $(C_h, D_h)$ has order $\le \frac{3}{4}\theta < \theta$, and $\mathscr{T}_j$ has order $\theta$, it follows that $(D_h, C_h) \in \mathscr{T}_j$. Since $\mathscr{L}_i$ $\theta$-isolates $\mathscr{T}_i$, we deduce that $A \subseteq C_h$ and $D_h \subseteq B$. Moreover, since $(C_h, D_h)$ distinguishes $\mathscr{T}_i$ from $\mathscr{T}_j$ and has order $\le \frac{3}{4}\theta$, it follows that the $(\mathscr{T}_i, \mathscr{T}_j)$-distinction $(A, B)$ also has order $\le \frac{3}{4}\theta$. This proves (1).

\hspace*{12pt}(2) \textit{Each member of $\mathscr{L}^*$ is titanic.}

\textbf{Subproof.} Let $(A, B) \in \mathscr{L}^*$, and choose $j$ as above. We claim that $(B, A)$ is linked to $\mathscr{T}_j$. For suppose that there exists $(B', A') \in \mathscr{T}_j$ of smaller order than $(A, B)$, and with $B \subseteq B'$ and $A' \subseteq A$. Since $(A, B) \in \mathscr{T}_i$ and $A' \subseteq A$ it follows that $(A', B') \in \mathscr{T}_i$, and so $(A', B')$ distinguishes $\mathscr{T}_i$ from $\mathscr{T}_j$; and hence has order at least that of $(A, B)$, a contradiction. Thus there is no such $(B', A')$, and so $(B, A)$ is linked to $\mathscr{T}_j$. Since $(B, A)$ has order $\le \frac{3}{4}\theta$ by (1), it follows that $(A, B)$ is titanic, by 3.3. This proves (2).

By (1), there exist $\mathscr{L}_1^*, \dots, \mathscr{L}_k^* \subseteq \mathscr{L}^*$, mutually disjoint and with union $\mathscr{L}^*$, such that for $1 \le g \le k$, every $(A, B) \in \mathscr{L}_g^*$ satisfies $A \subseteq C_g$ and $D_g \subseteq B$. Fix $h$ with $1 \le h \le k$.

\hspace*{12pt}(3) \textit{There is no tangle $\mathscr{T}$ of order $\theta$ in $G$ with $\mathscr{L}_h^* \cup \{(D_h, C_h)\} \subseteq \mathscr{T}$.}

\textbf{Subproof.} Suppose that $\mathscr{T}$ is such a tangle. From the hypothesis, there exists $j$ with $1 \le j \le n$ such that $\mathscr{L}_j \subseteq \mathscr{T}$. Since $(D_h, C_h) \in \mathscr{T}$ and $(C_h, D_h) \in \mathscr{L}_i$, it follows that $i \ne j$. Let $(A, B)$ be the $(\mathscr{T}_i, \mathscr{T}_j)$-distinction. Since $\mathscr{L}_j \nsubseteq \mathscr{T}_i$, there exists $(C, D) \in \mathscr{L}_j$ such that $(C, D) \notin \mathscr{T}_i$. Therefore $(D, C) \in \mathscr{T}_i$, since $\mathscr{L}_j$ has order $<\theta$, and hence $D \subseteq A$
